% problem_id: S001-lemma-coverings-site
% source_id: S001 (Stacks Project)
% source_file: sets.tex
% statement_locator: sets.tex lines 952--970
% proof_locator: sets.tex lines 972--1105
% env: lemma
% id_note: label 在本来源内唯一
% pairing: tight (verified)
% status: complete (statement + author proof verbatim + reason + locators)
%
\begin{lemma}
\label{lemma-coverings-site}
With notations as above.
Let $\text{Cov}_0 \subset \text{Cov}(\mathcal{C})$
be a set contained in $\text{Cov}(\mathcal{C})$.
There exist a cardinal $\kappa$ and a limit ordinal $\alpha$
with the following properties:
\begin{enumerate}
\item We have $\text{Cov}_0 \subset \text{Cov}(\mathcal{C})_{\kappa, \alpha}$.
\item The set of coverings
$\text{Cov}(\mathcal{C})_{\kappa, \alpha}$ satisfies
(1), (2), and (3) of Sites, Definition \ref{sites-definition-site} (see above).
In other words $(\mathcal{C}, \text{Cov}(\mathcal{C})_{\kappa, \alpha})$
is a site.
\item Every covering in $\text{Cov}(\mathcal{C})$
is combinatorially equivalent
to a covering in $\text{Cov}(\mathcal{C})_{\kappa, \alpha}$.
\end{enumerate}
\end{lemma}

\begin{proof}
To prove this, we first consider the set $\mathcal{S}$ of all
sets of morphisms of $\mathcal{C}$ with fixed target.
In other words, an element of $\mathcal{S}$ is a subset $T$
of $\text{Arrows}(\mathcal{C})$ such that all
elements of $T$ have the same target. Given a family
$\mathcal{U} = \{\varphi_i : U_i \to U\}_{i\in I}$ of morphisms with fixed
target, we define
$Supp(\mathcal{U}) = \{ \varphi \in \text{Arrows}(\mathcal{C})
\mid \exists i\in I, \varphi = \varphi_i\}$.
Note that two families $\mathcal{U} =  \{\varphi_i : U_i \to U\}_{i\in I}$
and $\mathcal{V} = \{V_j \to V\}_{j \in J}$ are combinatorially
equivalent if and only if $Supp(\mathcal{U}) = Supp(\mathcal{V})$.
Next, we define
$\mathcal{S}_\tau \subset \mathcal{S}$ to be the subset
$\mathcal{S}_\tau = \{ T \in \mathcal{S} \mid
\exists\ \mathcal{U} \in \text{Cov}(\mathcal{C}) \ T = Supp(\mathcal{U})\}$.
For every element $T \in \mathcal{S}_\tau$, set
$\beta(T)$ to equal the least ordinal $\beta$ such that
there exists a $\mathcal{U} \in \text{Cov}(\mathcal{C})_\beta$
such that $T = \text{Supp}(\mathcal{U})$. Finally, set
$\beta_0 = \sup_{T \in S_\tau} \beta(T)$.
At this point it follows that every $\mathcal{U} \in \text{Cov}(\mathcal{C})$
is combinatorially equivalent to some element
of $\text{Cov}(\mathcal{C})_{\beta_0}$.

\medskip\noindent
Let $\kappa$ be the maximum of $\aleph_0$,
the cardinality $|\text{Arrows}(\mathcal{C})|$,
$$
\sup\nolimits_{\{U_i \to U\}_{i\in I} \in \text{Cov}(\mathcal{C})_{\beta_0}}
|I|,
\quad\text{and}\quad
\sup\nolimits_{\{U_i \to U\}_{i\in I} \in \text{Cov}_0} |I|.
$$
Since $\kappa$ is an infinite cardinal, we have
$\kappa \otimes \kappa = \kappa$. Note that obviously
$\text{Cov}(\mathcal{C})_{\beta_0} =
\text{Cov}(\mathcal{C})_{\kappa, \beta_0}$.

\medskip\noindent
We define, by transfinite induction, a function $f$ which associates
to every ordinal an ordinal as follows. Let $f(0) = 0$.
Given $f(\alpha)$, we define $f(\alpha + 1)$ to be the least
ordinal $\beta$ such that the following hold:
\begin{enumerate}
\item We have $\alpha + 1 \leq \beta$ and $f(\alpha) \leq \beta$.
\item If $\{U_i \to U\}_{i\in I}
\in \text{Cov}(\mathcal{C})_{\kappa, f(\alpha)}$
and for each $i$ we have
$\{W_{ij} \to U_i\}_{j\in J_i}
\in \text{Cov}(\mathcal{C})_{\kappa, f(\alpha)}$,
then
$\{W_{ij} \to U\}_{i \in I, j\in J_i}
\in \text{Cov}(\mathcal{C})_{\kappa, \beta}$.
\item If $\{U_i \to U\}_{i\in I}
\in \text{Cov}(\mathcal{C})_{\kappa, \alpha}$
and $W \to U$ is a morphism of $\mathcal{C}$, then
$\{U_i \times_U W \to W \}_{i\in I}
\in \text{Cov}(\mathcal{C})_{\kappa, \beta}$.
\end{enumerate}
To see $\beta$ exists we note that clearly the collection of all
coverings $\{W_{ij} \to U\}$ and $\{U_i \times_U W \to W \}$ that occur in
(2) and (3) form a set. Hence there is some ordinal $\beta$ such that
$V_\beta$ contains all of these coverings. Moreover, the index set
of the covering $\{W_{ij} \to U\}$ has cardinality
$\sum_{i \in I} |J_i| \leq \kappa \otimes \kappa = \kappa$, and
hence these coverings are contained in
$\text{Cov}(\mathcal{C})_{\kappa, \beta}$.
Since every nonempty collection of ordinals has a least element
we see that $f(\alpha + 1)$ is well defined. Finally, if $\alpha$
is a limit ordinal, then we set
$f(\alpha) = \sup_{\alpha' < \alpha} f(\alpha')$.

\medskip\noindent
Pick an ordinal $\beta_1$ such that
$\text{Arrows}(\mathcal{C}) \subset V_{\beta_1}$,
$\text{Cov}_0 \subset V_{\beta_0}$,
and $\beta_1 \geq \beta_0$.
By construction $f(\beta_1) \geq \beta_1$ and we see that
the same properties hold for $V_{f(\beta_1)}$. Moreover, as $f$ is
nondecreasing this remains true for any $\beta \geq \beta_1$.
Next, choose any ordinal $\beta_2 > \beta_1$ with
cofinality $\text{cf}(\beta_2) > \kappa$. This is possible
since the cofinality of ordinals gets arbitrarily large, see
Proposition \ref{proposition-exist-ordinals-large-cofinality}.
We claim that the pair $\kappa$,
$\alpha = f(\beta_2)$ is a solution to the problem posed in the lemma.

\medskip\noindent
The first and third property of the lemma holds by our choices
of $\kappa$, $\beta_2 > \beta_1 > \beta_0$ above.

\medskip\noindent
Since $\beta_2$ is a limit ordinal (as its cofinality is infinite)
we get $f(\beta_2) = \sup_{\beta < \beta_2} f(\beta)$.
Hence $\{f(\beta) \mid \beta < \beta_2\} \subset f(\beta_2)$ is a
cofinal subset. Hence we see that
$$
V_\alpha = V_{f(\beta_2)} = \bigcup\nolimits_{\beta < \beta_2} V_{f(\beta)}.
$$
Now, let $\mathcal{U} \in \text{Cov}_{\kappa, \alpha}$.
We define $\beta(\mathcal{U})$ to be the least ordinal $\beta$ such that
$\mathcal{U} \in \text{Cov}_{\kappa, f(\beta)}$. By the above we see
that always $\beta(\mathcal{U}) < \beta_2$.

\medskip\noindent
We have to show properties (1), (2), and (3) defining a site
hold for the pair $(\mathcal{C}, \text{Cov}_{\kappa, \alpha})$.
The first holds because by our choice of $\beta_2$
all arrows of $\mathcal{C}$ are contained in $V_{f(\beta_2)}$.
For the third, we use that given a covering
$\mathcal{U} = \{U_i \to U\}_{i \in I}
\in \text{Cov}(\mathcal{C})_{\kappa, \alpha}$
we have $\beta(\mathcal{U}) < \beta_2$ and hence
any base change of $\mathcal{U}$ is by construction of
$f$ contained in $\text{Cov}(\mathcal{C})_{\kappa, f(\beta + 1)}$
and hence in $\text{Cov}(\mathcal{C})_{\kappa, \alpha}$.

\medskip\noindent
Finally, for the second condition, suppose that $\{U_i \to U\}_{i\in I}
\in \text{Cov}(\mathcal{C})_{\kappa, f(\alpha)}$
and for each $i$ we have
$\mathcal{W}_i = \{W_{ij} \to U_i\}_{j\in J_i}
\in \text{Cov}(\mathcal{C})_{\kappa, f(\alpha)}$.
Consider the function
$I \to \beta_2$, $i \mapsto \beta(\mathcal{W}_i)$. Since the cofinality
of $\beta_2$ is $> \kappa \geq |I|$ the image of this function cannot be a
cofinal subset. Hence there exists a $\beta < \beta_1$ such
that $\mathcal{W}_i \in \text{Cov}_{\kappa, f(\beta)}$ for all $i \in I$.
It follows that the covering $\{W_{ij} \to U\}_{i\in I, j \in J_i}$
is an element of $\text{Cov}(\mathcal{C})_{\kappa, f(\beta + 1)}
\subset \text{Cov}(\mathcal{C})_{\kappa, \alpha}$ as desired.
\end{proof}
